\begingroup
\setlength{\tabcolsep}{4.5pt}
\renewcommand{\arraystretch}{1.28}
\begin{longtable}{@{}>{\raggedright\arraybackslash}p{.14\textwidth}>{\raggedright\arraybackslash}p{.33\textwidth}>{\raggedright\arraybackslash}p{.25\textwidth}>{\raggedright\arraybackslash}p{.22\textwidth}@{}}
\caption{Réalisation physique des secteurs de jauge et scalaire.}\label{tab:masas-realizacion-gauge-escalar}\\
\toprule
\textbf{Classe} & \textbf{Objeto HMT} & \textbf{Application de réalisation} & \textbf{Conclusion et conditions} \\
\midrule
\endfirsthead
\toprule
\textbf{Classe} & \textbf{Objeto HMT} & \textbf{Application de réalisation} & \textbf{Conclusion et conditions} \\
\midrule
\endhead
\midrule
\multicolumn{4}{r}{\footnotesize Suite à la page suivante.}\\
\endfoot
\bottomrule
\endlastfoot
photon $\gamma$ & Carte nulle sectorielle et représentation de jauge correspondante. & Inclusion exacte de la section nulle dans la droite de masse. & Nullité exacte dans sa carte sectorielle. \\
gluones $g$ & Carte nulle sectorielle et représentation de jauge correspondante. & Inclusion exacte de la section nulle dans la droite de masse. & Nullité exacte dans sa carte sectorielle. \\
$W^\pm$ & parcours CT--108 W : 15/37, axe E-O, $(N,E,S,O)=(23,44,11,30)$, feuillet/résidu $(41,8)$, $(n_A,s_C)=(94,-1)$, $(W_+,W_-)=(563,565)$ ; caractère scalaire exact & Parcours CT--108 nominal $\to(n_A,s_C)\to(W_+,W_-)\to\chi_{\rm mass}\to\mathcal L_M$. & Parcours de jauge CT--108, signature et caractère scalaire déterminés avant la comparaison. \\
$Z^0$ & parcours CT--108 Z : 20/23, axe E-O, $(N,E,S,O)=(33,45,10,20)$, feuillet/résidu $(57,7)$, $(n_A,s_C)=(95,-1)$, $(W_+,W_-)=(569,571)$ ; caractère scalaire exact & Parcours CT--108 nominal $\to(n_A,s_C)\to(W_+,W_-)\to\chi_{\rm mass}\to\mathcal L_M$. & Parcours de jauge CT--108, signature et caractère scalaire déterminés avant la comparaison. \\
Higgs $H$ & parcours CT--108 H : 24/29, axe N-S, $(N,E,S,O)=(40,34,22,12)$, feuillet/résidu $(47,10)$, $(n_A,s_C)=(97,9)$, $(W_+,W_-)=(591,573)$ ; caractère scalaire exact & Parcours CT--108 nominal $\to(n_A,s_C)\to(W_+,W_-)\to\chi_{\rm mass}\to\mathcal L_M$. & Parcours scalaire CT--108, signature et caractère déterminés avant la comparaison. \\
\end{longtable}
\endgroup
